\chapter*{Abstrak}

% Abstrak berisikan resume yang menggambarkan keseluruhan rencana TA yang akan dikerjakan yang meliputi permasalahan, metodologi dan hipotesis awal. Secara umum poin penting yang harus ada didalam abstrak sebuah proposel TA adalah sebagai berikut:
% \begin{enumerate}
%     \item Deskripsi singkat permasalahan (1-2 kalimat)
%     \item Tujuan utama
%     \item Metode yang digunakan atau solusi yang ditawarkan (1-3 kalimat)
%     \item Rencana sumber dan/atau jenis data dan/atau studi kasus yang digunakan
%     \item Hipotesis awal
% \end{enumerate}

% Selain itu, pada abstrak harus dituliskan kata kunci atau keyword, yang berisikan kata-kata yang medeskripsikan isi tulisan dan ditulis dengan huruf non kapital. Kata kunci maksimum sebanyak 6 kata.
  
% \vspace{0.5 cm}
% \begin{flushleft}
% {\textbf{Kata Kunci:} keyword-1, keyword-2, keyword-3.}
%\end{flushleft}



Pengiriman \textit{invoice} memerlukan penentuan prioritas yang konsisten agar dokumen dapat diterima sesuai aturan \textit{cut-off} dan jadwal penerimaan pelanggan. Penelitian ini bertujuan melakukan \textit{comparative analysis} antara klasifikasi \textit{Rule-Based} yang dibangun dari SOP perusahaan dan \textit{Decision Tree} berbasis \textit{entropy} untuk menentukan prioritas pengiriman \textit{invoice}. Sistem pelacakan \textit{invoice} dan \textit{Proof of Delivery} (POD) ditempatkan sebagai konteks penerapan hasil klasifikasi.

Dataset analisis terdiri atas 99 \textit{invoice} unik dengan 30 kelas \textit{Urgent} dan 69 kelas \textit{Not Urgent}. Label historis admin perusahaan digunakan sebagai \textit{ground truth} bersama untuk melatih \textit{Decision Tree} dan mengevaluasi kedua metode. Evaluasi dilakukan melalui empat skenario, yaitu E1 \textit{stratified hold-out} 80:20, E2 \textit{stratified hold-out} 70:30, E3 \textit{stratified 5-fold cross-validation}, dan E4 \textit{leave-one-out cross-validation} (LOOCV).

Pada skenario LOOCV, \textit{Rule-Based} menghasilkan akurasi 93,94\%, sedangkan \textit{Decision Tree} menghasilkan akurasi 97,98\%. Namun, uji Exact McNemar menghasilkan \(p = 0{,}2891\). Hasil tersebut menunjukkan perbedaan performa secara deskriptif, tetapi belum menunjukkan perbedaan yang signifikan secara statistik. Temuan penelitian menjadi dasar pembahasan komparatif mengenai performa, pola kesalahan, dan karakteristik operasional kedua metode tanpa menyatakan keunggulan umum salah satu metode.

\vspace{0.5 cm} \begin{flushleft} {\textbf{Kata Kunci:}} \textit{invoice priority classification}, \textit{Rule-Based}, \textit{Decision Tree}, \textit{comparative analysis}, \textit{administrator ground truth}, LOOCV \end{flushleft}
